\documentclass[sigconf,nonacm]{acmart}

\usepackage{amsmath,amssymb}
\usepackage{url}

% Suppress ACM conference and copyright boilerplate to ensure strict 1-page fit
\settopmatter{printacmref=false}
\renewcommand\footnotetextcopyrightpermission[1]{}
\pagestyle{plain}

\title{PRISM: Evaluating Privacy-Constrained Customer Churn Classification}

% Compact author block to conserve vertical column space
\author{Abhishek Majumdar (UG), Dev Parsana (UG), Ravi Vashi (UG), Bharath Irukulapati (GR, Lead)}
\affiliation{%
  \institution{Department of Computer Science, University of Victoria}
  \city{Victoria}
  \country{Canada}
}

\begin{document}

\begin{abstract}
Predicting customer churn is vital for subscription services, but standard pipelines rely on continuous, invasive tracking of user activity. Under strict privacy regulations and cold-start conditions, collecting granular behavioral telemetry is often infeasible. We present \textbf{PRISM} (Privacy-Restricted Information for Subscriber Modeling) to investigate churn prediction under constrained feature visibility. We measure how performance degrades as behavioral features are removed, testing whether basic, non-intrusive account metadata can reliably identify at-risk subscribers.
Repository: \url{https://github.com/ibac03/PRISM-churn-privacy}.
\end{abstract}

\maketitle

\section{Introduction}
Customer retention is vital for subscription business models, as churn directly reduces recurring revenue. While early churn detection enables proactive retention offers, standard pipelines assume unrestricted access to fine-grained user telemetry, including app logs and interaction histories. However, strict privacy regulations (e.g., GDPR, CCPA) and user tracking opt-outs increasingly limit the behavioral data companies can store. This project investigates whether data mining models can accurately detect churn when restricted to non-intrusive billing and administrative attributes.

\section{Motivation and Significance}
\label{sec:motivation}
Retaining existing customers is significantly more cost-effective than acquiring new ones \cite{gupta2005valuing}. Detecting churn without invasive tracking provides two key benefits:
(1)~\textbf{Privacy Compliance \& Trust:} Minimizing data collection mitigates regulatory penalties and preserves consumer trust \cite{martin2017data}; and
(2)~\textbf{Cold-Start Feasibility:} New subscribers lack interaction history, rendering telemetry-heavy models ineffective during early contract phases.

\subsection{Motivating Example}
Consider a subscriber who opts out of app tracking. Because the provider possesses no activity logs, telemetry-dependent classifiers cannot evaluate this user. If a model can reliably predict churn using only administrative terms (e.g., month-to-month contract, electronic check payment, and tenure), the provider can intervene with targeted retention offers without invading user privacy.

\subsection{Gaps in Existing Solutions}
Most churn studies engineer dense features from usage logs to maximize benchmark scores \cite{verbeke2012building}, overlooking settings where telemetry is restricted. Consequently, there is limited understanding of the performance trade-offs of using a privacy-preserving subsets of data. PRISM addresses this gap by formalizing and benchmarking churn classification across restricted visibility tiers.

\section{Formal Problem Definition}
\label{sec:formal_problem}
Following the course taxonomy, customer churn prediction is formulated as a \textbf{Boolean Classification} task under a \textbf{Fully Supervised Learning} framework.

Let $O$ denote the set of customer observations, and let $C = \{0, 1\}$ represent the target classes, where $c = 1$ (Positive Class) indicates churn and $c = 0$ (Negative Class) indicates retention. The task is to learn a mapping function:
\begin{equation}
    f: O \rightarrow C \quad \text{such that} \quad f(o) = c, \quad \forall o \in O, \; c \in C
\end{equation}
where each $o \in O$ represents a customer described across restricted attribute spaces (from basic billing records to full service configurations).

\subsection{Running Example Instance}
Consider an observation $o \in O$ representing a customer account: $\langle\text{Tenure: 2 mos}$, $\text{Contract: Monthly}$, $\text{Payment: E-check}\rangle$. The learned function maps this observation to the positive class:
\begin{equation*}
    f(o) = 1 \quad (\text{Customer Churn})
\end{equation*}
This flags an at-risk subscriber based strictly on administrative settings, requiring no private usage telemetry.

\section{Team and Collaboration Justification}
Our team comprises three undergraduate students (A. Majumdar, D. Parsana, R. Vashi) and one graduate student acting as project lead (B. Irukulapati). We formed this group around a shared interest in privacy-aware machine learning. Our collective skills across software development, data pipelines, and machine learning complement each other, enabling balanced collaboration on implementation, experimentation, and milestone deliverables.

% Compressed bibliography to guarantee clean termination on Page 1
\begin{thebibliography}{3}
\small
\setlength{\itemsep}{0pt}
\bibitem{gupta2005valuing}
S.~Gupta et~al. 2006. Modeling customer lifetime value. \emph{Journal of Service Research} 9, 2 (2006), 139--155.

\bibitem{martin2017data}
K.~Martin and N.~Murphy. 2017. The role of data privacy in marketing. \emph{J. Acad. Mark. Sci.} 45, 2 (2017), 135--155.

\bibitem{verbeke2012building}
W.~Verbeke et~al. 2012. New insights into churn prediction in the telecommunication sector. \emph{EJOR} 218, 1 (2012), 211--229.
\end{thebibliography}

\newpage
\appendix
\section{AI-Assisted Editing Disclosure (Category C)}
In accordance with course guidelines for Category C (AI-Assisted Editing), generative AI was used solely to assist with LaTeX formatting, syntax adjustments, and layout condensing to satisfy the strict single-page limit. All problem formulations, conceptual ideas, team details, and definitions were developed by the team members.

\subsection{Original Unedited Student Draft}
Please refer to the original draft document submitted along with this document.

\end{document}